\documentclass[12pt,a4paper]{article}

% ------------------- Пакеты -------------------
\usepackage[utf8]{inputenc}    % Кодировка
\usepackage[russian]{babel}    % Русский язык
\usepackage{graphicx}           % Для вставки картинок
\usepackage{hyperref}           % Гиперссылки
\usepackage{geometry}           % Настройка полей
\geometry{top=2cm, bottom=2cm, left=2.5cm, right=1.5cm}
\usepackage{titlesec}           % Настройка заголовков
\usepackage{listings}           % Оформление листингов кода
\usepackage{xcolor}              % Цвета для кода
\usepackage{enumitem}           % Улучшенные списки
\usepackage{fancyhdr}           % Колонтитулы
\usepackage{lastpage}           % Номер последней страницы
\usepackage{tcolorbox}           % Цветные блоки (для предупреждений)

% ------------------- Настройка кода -------------------
\lstdefinestyle{codestyle}{
    backgroundcolor=\color{lightgray!10},
    commentstyle=\color{green!60!black},
    keywordstyle=\color{blue},
    numberstyle=\tiny\color{gray},
    stringstyle=\color{orange},
    basicstyle=\ttfamily\footnotesize,
    breakatwhitespace=false,
    breaklines=true,
    captionpos=b,
    keepspaces=true,
    numbers=left,
    numbersep=5pt,
    showspaces=false,
    showstringspaces=false,
    showtabs=false,
    tabsize=2,
    language=Python,  % Можно заменить на C++, Bash и т.д.
    frame=single
}
\lstset{style=codestyle}

% ------------------- Настройка заголовков -------------------
\titleformat{\section}{\large\bfseries}{Пункт \thesection}{1em}{}
\titleformat{\subsection}{\normalsize\bfseries}{\thesubsection}{1em}{}

% ------------------- Настройка колонтитулов -------------------
\pagestyle{fancy}
\fancyhf{}
\fancyhead[L]{\leftmark}
\fancyhead[R]{\thepage/\pageref{LastPage}}
\renewcommand{\headrulewidth}{0.4pt}

% ------------------- Титульная информация -------------------
\title{\textbf{Технический мануал: \\ Работа с TOF-сенсором и UART}}
\author{Иван Иванов \\ ООО "ТехноПроект"}
\date{\today}

\begin{document}

% ------------------- Титульный лист -------------------
\begin{titlepage}
    \begin{center}
        \vspace*{1cm}
        \Huge
        \textbf{Информационное пособие к проекту по информатике}
        
        \vspace{0.5cm}
        \LARGE
        Московский физико-технический институт.\\
        Ребенков Андрей, Троянова Милана, Шапиев Халид.\\
        Группа Б03-502
        
        \vspace{1.5cm}
        \textbf{TOF-сенсор, UART интерфейс, COM-порты, обработка данных,
        GUI}
        
        \vfill
        \includegraphics[width=0.4\textwidth]{example-image} % Замените на логотип
        
        \vspace{0.8cm}
        \Large
        Долгопрудный, 2026
    \end{center}
\end{titlepage}

% ------------------- Оглавление -------------------
\tableofcontents
\newpage

% ------------------- Введение -------------------
\section{TOF-сенсор, кто это?}
Time-of-flight-камера (англ. Time-of-flight camera (ToF camera)) — видеокамера, формирующая так называемое дальностное изображение (дальностный портрет). Используется для создания изображений, которые в качестве пикселей содержат оценки расстояний от экрана до конкретных точек наблюдения.\\

\subsection{Задачи по сенсору}
\begin{itemize}
    \item создание обработчика данных(см. далее)
    \item Настройка UART-соединения.
    \item Разработка практичного и отвечающего требованиям приложения с GUI для работы с сенсором.
    \item Создание ПО для соединения сенсора и приложения.
    \item Возможны дополнения\\

    *обрабочик необходим для превращения сырых данных сенсора в шестнадцатиричной системе в что-то по типу: карты глубины, координатного файла с отметкой глубины, 3д модели глубины 2д изображения.
\end{itemize}

\subsection{Команда проекта}
\begin{description}[style=nextline]
    \item[Андрей] Специалист в создании красивых и функциональных оболочек для ПО. Фанат вархаммера.
    \item[Милана] Активный деятель команды, я хз что написать.
    \item[Халид] Лысый прогер, using namespace std не про него.
\end{description}

\newpage

% ------------------- Основная часть -------------------
\section{UART-интерфейс. Важные моменты.}

UART (Universal Asynchronous Receiver-Transmitter) - интерфейс связи между цифровыми устройствами, по нему МК может общаться с другими микросхемами и МК. В нашей работе это главное, для коннекта между сенсором и ПО.\\
UART не является шиной (но может использоваться как шина, топология выстраивается программистом) - он предназначен для передачи данных между двумя устройствами, причём оба этих устройства могут передавать данные по своему усмотрению, в том числе одновременно. Таким образом, UART - это просто передатчик между равнозначными "устройством 1" и "устройством 2".\\
*Шина (Bus) — это общий физический канал (набор проводников), который объединяет несколько устройств для передачи данных.\\
\subsection{Схема соединений}
\begin{itemize}
    \item \textbf{Сенсор (TX)} $\rightarrow$ \textbf{Переходник (RX)}
    \item \textbf{Сенсор (RX)} $\rightarrow$ \textbf{Переходник (TX)}
    \item \textbf{Сенсор (GND)} $\rightarrow$ \textbf{Переходник (GND)}
    \item \textbf{Сенсор (VCC)} $\rightarrow$ \textbf{3.3V / 5V (в зависимости от модели)}
\end{itemize}

TX - (transmit) пин для отправки данных, а RX - (receiver) получает их. То, как у нас будут соединен интерфейс с сенсором видно выше.

\begin{center}
    \fbox{Схема подключения (вставьте изображение здесь)}
\end{center}

\begin{tcolorbox}[colback=red!5!white, colframe=red!75!black, title=Важно!]
Перед подключением убедитесь, что уровни напряжений сенсора и переходника совпадают. Использование 5V логики для 3.3V сенсора может вывести его из строя.
\end{tcolorbox}

\subsection{Источники для ознакомления с UART}
\begin{itemize}
    \item \href{https://components.espressif.com/components/plasmapper/pl_uart/versions/1.1.0/examples?language=en}{интересная C++ обвязка для UART от пользователя plasmapper. Примеры демонстрируют инициализацию и даже создание эхо-сервера с использованием классов — то, что нужно для объектно-ориентированного подхода на микроконтроллере.}
    \item \href{https://doc.embedfire.com/linux/rk356x/Qt/zh/latest/lubancat_qt/senior/qt_serialport.html}{Мануал к библиотеке Qt- возможно ключевой для нашей работы с UART интерфейсом сенсора. Оригинал на китайском, но переводчик в помощь.}
\end{itemize}

\section{Настройка UART в Windows}

После подключения USB-UART переходника в системе будет создан виртуальный COM-порт. В отличие от Linux, где порт выглядит как файл \texttt{/dev/ttyUSB0}, в Windows он получает имя \texttt{COM3}, \texttt{COM5} и т.д.

\subsection{Проверка подключения (определение номера порта)}
\begin{enumerate}
    \item Откройте \textbf{Диспетчер устройств} (Device Manager). Самый быстрый способ: нажмите \texttt{Win + X} и выберите соответствующий пункт меню.
    \item Разверните раздел \textbf{Порты (COM и LPT)}.
    \item Найдите ваш переходник. Обычно он подписан названием чипа: \texttt{USB Serial Port (COM3)} или \texttt{CH340 (COM5)}.
    \item Запомните номер порта (например, \texttt{COM3}). Он понадобится для подключения.
\end{enumerate}

\begin{tcolorbox}[colback=blue!5!white, colframe=blue!75!black, title=Примечание]
Если переходник определился как неизвестное устройство (с желтым восклицательным знаком), необходимо установить драйвер. Для популярных чипов (CH340, CP2102, FT232) драйверы обычно скачиваются с сайта производителя или находятся автоматически через Windows Update.
\end{tcolorbox}

\subsection{Установка прав доступа}
В Windows, в отличие от Ubuntu, нет ограничений по группам для доступа к COM-портам. Если порт отображается в Диспетчере устройств без ошибок, то любая программа от имени текущего пользователя сможет к нему подключиться.
\begin{itemize}
    \item Единственная проблема: порт может быть занят другой программой.
    \item Решение: закройте все программы, которые могли открыть порт (среды разработки, предыдущие сессии терминала).
\end{itemize}

\subsection{Проверка приема данных (программа на Qt)}

Вместо готовых терминалов (PuTTY, Tera Term) мы будем использовать программу, написанную на Qt с использованием модуля \texttt{QSerialPort}. Это позволит нам не только проверить работу сенсора, но и получить основу для дальнейшей разработки проекта.

\subsubsection*{Создание простейшего терминала в Qt}

\begin{enumerate}
    \item \textbf{Создайте новый проект} в Qt Creator: 
    \begin{itemize}
        \item \texttt{File → New File or Project → Application → Qt Widgets Application}.
        \item Назовите проект, например, \texttt{TOF\_Sensor\_Test}.
    \end{itemize}
    
    \item \textbf{Добавьте модуль последовательного порта} в файл \texttt{TOF\_Sensor\_Test.pro}:
    \begin{lstlisting}[style=myCStyle]
QT += core gui serialport
    \end{lstlisting}
    
    \item \textbf{Разработайте интерфейс} (файл \texttt{mainwindow.ui}) или создайте его программно. Минимально необходимые элементы:
    \begin{itemize}
        \item Выпадающий список (\texttt{QComboBox}) для выбора COM-порта.
        \item Поле ввода (\texttt{QLineEdit}) или выпадающий список для выбора скорости (baudrate).
        \item Кнопка (\texttt{QPushButton}) "Подключиться" / "Отключиться".
        \item Текстовая область (\texttt{QTextEdit}) для вывода полученных данных.
    \end{itemize}
    
    \item \textbf{Реализуйте логику работы} в файле \texttt{mainwindow.cpp}:
\end{enumerate}

\begin{lstlisting}[style=myCStyle, caption=Пример кода для проверки приема данных]
#include "mainwindow.h"
#include "ui_mainwindow.h"
#include <QSerialPort>
#include <QSerialPortInfo>
#include <QDebug>

MainWindow::MainWindow(QWidget *parent) :
    QMainWindow(parent),
    ui(new Ui::MainWindow)
{
    ui->setupUi(this);
    
    // Поиск доступных портов при запуске
    foreach (const QSerialPortInfo &portInfo, QSerialPortInfo::availablePorts()) {
        ui->comboBoxPort->addItem(portInfo.portName());
    }
    
    // Настройка скоростей (baudrates)
    ui->comboBoxBaud->addItems(QStringList() << "9600" << "19200" << "38400" 
                                             << "57600" << "115200");
    ui->comboBoxBaud->setCurrentText("115200"); // По умолчанию
    
    // Создание объекта последовательного порта
    serial = new QSerialPort(this);
    
    // Связываем сигнал готовности данных с обработчиком
    connect(serial, &QSerialPort::readyRead, this, &MainWindow::readData);
    
    // Кнопка подключения
    connect(ui->pushButtonConnect, &QPushButton::clicked, this, &MainWindow::onConnectButtonClicked);
}

MainWindow::~MainWindow()
{
    if(serial->isOpen()) serial->close();
    delete ui;
}

void MainWindow::onConnectButtonClicked()
{
    if(serial->isOpen()) {
        // Отключаемся
        serial->close();
        ui->pushButtonConnect->setText("Подключиться");
        ui->textEdit->append("Порт закрыт");
    } else {
        // Настраиваем и подключаемся
        serial->setPortName(ui->comboBoxPort->currentText());
        serial->setBaudRate(ui->comboBoxBaud->currentText().toInt());
        serial->setDataBits(QSerialPort::Data8);
        serial->setParity(QSerialPort::NoParity);
        serial->setStopBits(QSerialPort::OneStop);
        serial->setFlowControl(QSerialPort::NoFlowControl);
        
        if(serial->open(QIODevice::ReadWrite)) {
            ui->pushButtonConnect->setText("Отключиться");
            ui->textEdit->append("Подключено к " + serial->portName() + 
                                 " со скоростью " + ui->comboBoxBaud->currentText());
        } else {
            ui->textEdit->append("Ошибка открытия порта: " + serial->errorString());
        }
    }
}

void MainWindow::readData()
{
    // Чтение всех доступных данных
    QByteArray data = serial->readAll();
    
    // Преобразование в hex-строку (для отладки)
    QString hexString = data.toHex(' ').toUpper();
    
    // Вывод в текстовое поле
    ui->textEdit->append("Получено (" + QString::number(data.size()) + " байт): " + hexString);
    
    // Дополнительно: попытка интерпретировать как ASCII, если это текст
    QString asciiString = QString::fromUtf8(data);
    if(!asciiString.isEmpty() && asciiString.isPrint()) {
        ui->textEdit->append("Как текст: " + asciiString);
    }
}
\end{lstlisting}

\begin{center}
    \fbox{Скриншот окна программы на Qt (вставьте изображение)}
\end{center}

\subsubsection*{Запуск и проверка}

\begin{enumerate}
    \item Подключите ваш TOF-сенсор к компьютеру через USB-UART переходник.
    \item Запустите программу (зеленая кнопка \textbf{Run} в Qt Creator или собранный \texttt{.exe} файл).
    \item Выберите нужный порт (например, \texttt{COM3}) в выпадающем списке.
    \item Установите скорость, соответствующую даташиту вашего сенсора (обычно 115200).
    \item Нажмите кнопку \textbf{Подключиться}.
    \item Если сенсор отправляет данные, они начнут появляться в текстовой области в 16-ричном виде.
\end{enumerate}

\begin{tcolorbox}[colback=green!5!white, colframe=green!75!black, title=Преимущества подхода на Qt]
\begin{itemize}
    \item \textbf{Полный контроль:} Вы сами решаете, как отображать и обрабатывать данные.
    \item \textbf{Обработка hex:} Функция \texttt{data.toHex()} делает именно то, что нужно для вашего сенсора.
    \item \textbf{Отладка:} Код можно легко дополнить логированием в файл, построением графиков и т.д.
    \item \textbf{Кроссплатформенность:} Этот же код будет работать на Linux и macOS без изменений.
    \item \textbf{Нет проблем с "зависшими" портами:} Приложение само корректно закрывает порт при выходе.
\end{itemize}
\end{tcolorbox}

\begin{tcolorbox}[colback=red!5!white, colframe=red!75!black, title=Важно!]
\begin{itemize}
    \item Убедитесь, что в \texttt{.pro} файле проекта добавлена строка \texttt{QT += serialport}.
    \item Если данные не отображаются, проверьте:
    \begin{itemize}
        \item Правильность выбранного COM-порта (в Диспетчере устройств Windows).
        \item Соответствие скорости (baudrate) даташиту сенсора.
        \item Не запущена ли другая программа (Arduino IDE, монитор порта), которая могла захватить порт.
    \end{itemize}
    \item При закрытии программы порт освобождается автоматически.
\end{itemize}
\end{tcolorbox}

\section{Программная обработка данных}

Данные от TOF-сенсора обычно приходят в бинарном виде. Разберем пример на Python с использованием библиотеки \texttt{pyserial}.

\subsection{Установка библиотеки}
\begin{lstlisting}[language=bash]
pip install pyserial
\end{lstlisting}

\subsection{Пример чтения сырых данных}
\begin{lstlisting}[language=Python]
import serial

# Открытие порта
ser = serial.Serial(
    port='/dev/ttyUSB0',
    baudrate=115200,
    timeout=1,
    bytesize=8,
    parity='N',
    stopbits=1
)

while True:
    # Чтение 9 байт (для сенсоров типа TF-Luna)
    data = ser.read(9)
    if data and len(data) == 9:
        print(f"Получено байт: {data.hex()}")
\end{lstlisting}

\subsection{Парсинг пакета}
Предположим, что сенсор отправляет пакеты формата:

\begin{center}
\begin{tabular}{|c|c|c|c|c|c|c|c|c|}
\hline
0x59 & 0x59 & DIST\_L & DIST\_H & STR\_L & STR\_H & TEMP\_L & TEMP\_H & SUM \\
\hline
\end{tabular}
\end{center}

Где:
\begin{itemize}
    \item \texttt{DIST\_L} и \texttt{DIST\_H} — младший и старший байты дистанции.
    \item \texttt{SUM} — контрольная сумма (сумма первых 8 байт \& 0xFF).
\end{itemize}

Функция парсинга:
\begin{lstlisting}[language=Python]
def parse_tf_packet(packet):
    # Проверка заголовка
    if packet[0] != 0x59 or packet[1] != 0x59:
        return None
    
    # Контрольная сумма
    checksum = sum(packet[:8]) & 0xFF
    if checksum != packet[8]:
        return None
    
    # Дистанция (мм)
    distance = packet[2] + packet[3] * 256
    
    # Сила сигнала
    strength = packet[4] + packet[5] * 256
    
    # Температура (специфично для TF-Luna)
    temp = packet[6] + packet[7] * 256
    temp_c = temp / 8.0 - 256
    
    return distance, strength, temp_c
\end{lstlisting}

\subsection{Цикл чтения с буферизацией}
Чтобы не потерять данные при обрыве пакета, используем буфер:
\begin{lstlisting}[language=Python]
buffer = bytearray()

while True:
    # Добавляем новые данные в буфер
    buffer += ser.read(ser.in_waiting)
    
    # Ищем заголовок
    while len(buffer) >= 9:
        if buffer[0] == 0x59 and buffer[1] == 0x59:
            packet = buffer[:9]
            result = parse_tf_packet(packet)
            if result:
                dist, strength, temp = result
                print(f"Дистанция: {dist} мм, Сила: {strength}, Темп: {temp:.1f}C")
                buffer = buffer[9:]  # Удаляем обработанное
            else:
                buffer.pop(0)  # Сдвиг на байт
        else:
            buffer.pop(0)
\end{lstlisting}

\newpage
\section{Возможные проблемы и их решение}

\begin{enumerate}
    \item \textbf{Порт не открывается: Permission denied.}
        \begin{quote}
            Решение: Добавьте пользователя в группу dialout.
        \end{quote}
    
    \item \textbf{Вижу мусор или данные не приходят.}
        \begin{quote}
            Проверьте скорость (baudrate) в коде и даташите. Для TF-Luna часто используется 115200, но может быть и 9600.
        \end{quote}
    
    \item \textbf{Пакеты "рвутся", расстояние скачет.}
        \begin{quote}
            Используйте буферизацию и проверку контрольной суммы (как в примере выше).
        \end{quote}
\end{enumerate}

\begin{tcolorbox}[colback=yellow!5!white, colframe=yellow!75!black, title=Рекомендация]
При работе с UART всегда добавляйте таймаут (timeout) при открытии порта, чтобы программа не зависала, если сенсор отключится.
\end{tcolorbox}

\section{Приложение А. Справочная информация}

\subsection{Часто используемые скорости UART}
\begin{itemize}
    \item 9600 бод
    \item 19200 бод
    \item 38400 бод
    \item 57600 бод
    \item 115200 бод (наиболее распространенная для TOF)
\end{itemize}

\section{GUI/Интерфейс приложения}

\hyperlink{https://doc.qt.io/}{Возможно}, только это и нужно.

Qt должен отвечать всем нашим требованиям в разработке приложения. Вот примеры программ, которые были сделаны через него:

\begin{itemize}
    \item KDE Plasma
    \item Telegram desktop
    \item Virtualbox
\end{itemize}

Так что наверное в нем и будем работать.
% ------------------- Заключение -------------------
\section{Заключение}
В данном мануале рассмотрены основные этапы подключения TOF-сенсора к ПК и обработки данных. Предложенные примеры кода могут быть адаптированы под конкретную модель сенсора.

\end{document}